\documentclass[10pt,a4paper]{amsart}
\usepackage[margin=19mm]{geometry}
\usepackage{amsmath,amssymb,mathtools}
\usepackage[scr=rsfs,cal=euler]{mathalfa}
\usepackage{xfrac}
\usepackage[unicode,hidelinks,bookmarks=false]{hyperref}
\newcommand{\ie}{i.e.\ }
\newcommand{\cf}{cf.\ }
\newcommand{\resp}{respectively\ }
\newcommand{\Subsection}{\subsection}
\providecommand{\nobreakdash}{\nobreak\hbox{-}}
\setlength{\emergencystretch}{2em}
\multlinegap0pt
\usepackage{bm}
\usepackage{bbm}
\usepackage{dsfont}
\usepackage{xspace}
\usepackage{cancel}
\usepackage{mathtools}
\usepackage{amsthm}
\makeatletter
\@ifundefined{proposition}{\newtheorem{proposition}{Proposition}}{}
\@ifundefined{theorem}{\newtheorem{theorem}{Theorem}}{}
\makeatother
\def\A{\mathcal A}
\title[Enumerative Combinatorics of Homogeneous Linear Orderings]{Enumerative Combinatorics of Homogeneous Linear Orderings}
\author{David Gonzalez}
\date{Source: arXiv:2604.14255v1 (CC BY 4.0)}
\begin{document}
\maketitle

\noindent\emph{Source and licence.} Enumerative Combinatorics of Homogeneous Linear Orderings. Version arXiv:2604.14255v1 (CC BY 4.0), distributed under the Creative Commons Attribution 4.0 International licence, as recorded in the arXiv licence metadata for this exact version and shown on the abstract page https://arxiv.org/abs/2604.14255v1. Full rights notice: licenses/arxiv-2604.14255-CC-BY-4.0.txt.\par\smallskip

\noindent\textit{Editorial scope.} Counted unit: the Theorem whose source label is \texttt{None} in the article (source0.tex lines 611--629, source label \texttt{None}), delivered together with the article's own complete original proof of it (source0.tex lines 631--699, one \texttt{proof} environment of 69 lines). No mathematics is written, repaired or re-derived: statement and proof are the article's own text line for line. Required context retained uncounted: prop:coloriffSp (lines 349--351); prop:spImage (lines 458--460); thm:CnmClassify (lines 590--597). Every definition, display and label of those lines is retained unchanged. Omitted from this package and not redistributed: 1--348, 352--457, 461--589, 598--610, 630--630, 700--898. Nothing inside a retained excerpt is omitted or altered: every retained body is delivered exactly as the article's lines are written, so no omission receipt of any kind accompanies this package. Citations: the retained excerpts carry no citation command at all, so no bibliography entry of the article is transcribed and no reference is redistributed. Reference adaptation: none was needed and none was made; every reference inside the delivered unit resolves inside it, so no unresolved reference or citation is produced. Printed slips kept literal: no printed word, symbol, hypothesis or number of the counted statement or of its proof was altered. Layout and class: the article's own class is \texttt{amsart} with packages including amssymb, xy, hyperref, color, pdfsync, graphicx, eucal, lmodern, cleveref, amsmath, amsthm, mathtools, MnSymbol, setspace; the wrapper is the lane's pinned \texttt{amsart} editorial wrapper at 10pt on A4, which restates the theorem-like environments the retained text needs and adds the article's own macro definitions that the retained text needs (\texttt{\textbackslash A}), with extra wrapper packages (bm, bbm, dsfont, xspace, cancel, mathtools). The delivered numbering is the unit's own. Assets: no retained span contains a figure environment, a graphics-inclusion command, a PDF-page inclusion, a table or a TikZ picture; no image file is copied into this package, the package declares no asset dependency and no image file is redistributed with it. Licence: version arXiv:2604.14255v1 of this article is distributed under the Creative Commons Attribution 4.0 International licence, as recorded in the arXiv licence metadata for this exact version and shown on the abstract page https://arxiv.org/abs/2604.14255v1. A full rights notice is delivered at licenses/arxiv-2604.14255-CC-BY-4.0.txt. Characters: every retained excerpt is pure ASCII, so the delivered engine, XeLaTeX with its default Latin Modern text font, reports no missing character.
\begin{proposition}\label{prop:coloriffSp}
There are inverse embeddings $\Phi:Mod(S)\to Mod(T)$ and $\Psi:Mod(T)\to Mod(S)$.
\end{proposition}

\begin{proposition}\label{prop:spImage}
    There is a computable injective map $F:Mod(T')\to Mod(LO)$ such that the image of $F$ is exactly the sp-homogeneous linear orderings. 
\end{proposition}

\begin{theorem}\label{thm:CnmClassify}
    \begin{enumerate}
        \item If $n,m<\infty$ then the $C_{n,m}$-homogeneous structures are the sp-homogeneous structures with blocks of size at most $n+m+1$.
        \item If $n=\infty$ and $m<\infty$ then the $C_{n,m}$-homogeneous structures are the sp-homogeneous structures without blocks isomorphic to $\omega$.
        \item If $m=\infty$ and $n<\infty$ then the $C_{n,m}$-homogeneous structures are the sp-homogeneous structures without blocks isomorphic to $\omega^*$.
        \item If $m=n=\infty$ then the $C_{n,m}$-homogeneous structures are exactly the sp-homogeneous structures.
    \end{enumerate}
\end{theorem}

\begin{theorem}
Let $n$ and $m$ be finite.
\begin{enumerate}
    \item There are only finitely many $C_{n,m}$-homogeneous linear orderings.
    \item The number of $C_{n,m}$-homogeneous linear orderings only depends on $k=n+m+1$.
    \item If $I(k)$ is the number of $C_{n,m}$-homogeneous linear orderings with $k=n+m+1$, then $I(k)$ satisfies the following recursive definition.
    Let
    \[K_1(0)=1,K_2(0)=0 \text{ and }\] 
    \[K_1(k+1)= \sum_{i=0}^{k} {k+1 \choose i} \big(K_1(i)+K_2(i)\big) \text{ and } K_2(k+1)=(k+1)K_1(k).\]
    Then
    \[I(k)=\sum_{i=0}^{k} \big(K_1(i)+K_2(i)\big) {k \choose i}.\]
    \item $I(k)$ has the closed form
    \[I(k)=\sum_{m=1}^k{k \choose m}\sum_{n=1}^m\sum_{r=0}^{\lceil \frac{n}2 \rceil}{n-r+1 \choose r}{m \choose r}r!(n-r)!S(m-r,n-r),\]
    Where $S(n,m)$ are the Stirling numbers of the second kind.
    \item $I(k)$ is equal to the number of homogeneous linear orderings in $k$ colors satisfying $T$.
    \item $I(k)$ is equal to the number of models of $T'_k$, where $T'_k$ is the reduction of the theory $T'$ to the language $R_1,\dots,R_k$ and $S_1,\dots,S_k$.
\end{enumerate}
     
\end{theorem}

\begin{proof}
    Note that (2) follows immediately from Theorem \ref{thm:CnmClassify} and that (1) will follow from (3).
    By Proposition \ref{prop:spImage} and Theorem \ref{thm:CnmClassify}, every $C_{n,m}$-homogeneous linear ordering lies in the image $F(T')$.
    Furthermore, following the definition of $F$, any element with predicate outside of $R_1,\dots,R_k$ and $S_1,\dots,S_k$ results in a block of size at least $k+1$ in $F(T')$.
    Conversely, if only $R_1,\dots,R_k$ and $S_1,\dots,S_k$ are used, only blocks of size at most $k$ are produced in $F(T')$.
    This means that $F$ is a bijection between the models of $T'_k$ and the $C_{n,m}$-homogeneous linear orderings.
    The same analysis for $\Psi$ from Proposition \ref{prop:coloriffSp} shows that the homogeneous linear orderings satisfying $T$ in $k$ colors (only the colors $c_1,\dots,c_k)$ are in bijection with the $C_{n,m}$-homogeneous linear orderings.
    In particular, (5) and (6) follow at once.

    All that remains is demonstrating (3) and (4).
    To accomplish this, we will explicitly count the models of $T'_k$, and appeal to the above analysis to observe that this gives the other desired counts as well.
    We begin with (3).
    We interpret $K_1(i)$ as the number of models of $T'_i$ using all $i$ colors where the first element satisfies some $R_i$, and interpret $K_2(i)$ as the number of models of $T'_i$ using all $i$ colors where the first element satisfies some $S_i$.
    Note that $\A\models T'_{k+1}$ must have finitely many elements; no color can be reused, and every element takes at least one color.
    Therefore, there is always a first element.
    If we are able to demonstrate that $K_1(i)$ and $K_2(i)$ act as we described above, the formula
    \[I(k)=\sum_{i=1}^{k} \big(K_1(i)+K_2(i)\big) {k \choose i},\]
    follows at once.
    This is because a model of $T'_k$ uses $i\leq k$ colors.
    It has ${k \choose i}$ options for that set of colors and $K_1(i)+K_2(i)$ many underlying models for that choice of colors.
    We focus on demonstrating that $K_1$ and $K_2$ are as desired.
    We demonstrate this fact inductively.
    Say that this characterization holds for $i\leq k$.
    Consider a model $\A\models T'_{k+1}$ using all of the colors.
    Say that the first element of $\A$ has $k+1-i$ of the $S_j$ hold of it.
    Note that there are ${k+1 \choose i}$ choices for which of the colors are used.
    The remaining ordering is a model of $T_i'$, of which there are $K_1(i)+K_2(i)$ many by induction.
    Therefore,
    \[K_1(k+1)= \sum_{i=0}^{k} {k+1 \choose i} \big(K_1(i)+K_2(i)\big),\]
    as desired.
    Now consider the case that the first element of $\A$ has some $R_j$ hold of it.
    Note that there are $k+1$ choices for which of the colors is used.
    The remaining ordering is a model of $T_k'$, but it necessarily does not begin with an element with some $R_i$ holding (otherwise, two such points would be adjacent, which is forbidden by our axioms).
    There are $K_1(k)$ many of these models by induction.
    Therefore,
    \[K_2(k+1)=(k+1)K_1(k),\]
    as desired.

    We now demonstrate (4) by counting the models of $T_k'$ directly.
    Each term in the sum
    \[I(k)=\sum_{m=1}^k{k \choose m}\sum_{n=1}^m\sum_{r=0}^{\lceil \frac{n}2 \rceil}{n-r+1 \choose r}{m \choose r}r!(n-r)!S(m-r,n-r),\]
    will be interpreted explicitly.
    Interpret $m$ as the total number of colors used in $T_k'$, and ${k \choose m}$ counts the number of choices for said colors.
    Therefore, it is enough to show that there are
    \[\sum_{n=1}^m\sum_{r=0}^{\lceil \frac{n}2 \rceil}{n-r+1 \choose r}{m \choose r}r!(n-r)!S(m-r,n-r)\]
    many models of $T_m'$ using all of the $m$ colors.
    Interpret $n$ as the number of elements in the model.
    Interpret $r$ as the number of elements satisfying some $R_i$ in the model.
    As these elements cannot be successive, there are at most $\lceil \frac{n}2 \rceil$ such elements.
    The isomorphism type of the model is now determined by
    \begin{enumerate}
        \item A choice of which elements are $R_i$ elements
        \item A choice of the colors used on the $R_i$ elements
        \item A choice of the order of the colors on the $R_i$ elements
        \item A choice of the order of the color sets for the $S_i$ elements
        \item A partition of the colors among the $S_i$ elements.
    \end{enumerate}
    Note that each of these choices is made independently.
    Therefore, we can count each quantity and multiply them together to get the final number of models subject to the parameters $k,m,n,r$.
    The counts are as follows:
    \begin{enumerate}
        \item ${n-r+1 \choose r}$, which is a known formula for the number of length $n$ binary sequences with $r$ ones, no two of which are consecutive.
        \item ${m \choose r}$, by definition
        \item $r!$, by definition
        \item $(n-r)!$, by definition
        \item $S(m-r,n-r)$, as we are splitting $m-r$ ordered colors among $n-r$ ordered points, and the Stirling numbers of the second kind count exactly this quantity.
    \end{enumerate}
    This completes the proof.
\end{proof}

\end{document}
